\subsection{PH3 implies PH2 for the representation-theoretic stretching families}

\begin{proposition}[P-F12: the denominator length is controlled by Weyl growth]
\label{prop:ph3ph2}
Let $G,H$ be connected reductive complex algebraic groups, let
$\rho:H\to G$ be an algebraic homomorphism, and let $\lambda,\pi$ be
dominant integral weights. Put
\[
 f(n)=\dim\operatorname{Hom}_H
       (V_{n\pi}(H),V_{n\lambda}(G)|_H).
\]
Suppose this is a nonzero global quasi-polynomial on $n\ge0$ and has a
positive rational form
\[
 \sum_{n\ge0}f(n)T^n
 =\frac{N(T)}{\prod_{j=1}^{r}(1-T^{a_j})},
 \quad N\in\mathbb Z_{\ge0}[T]\setminus\{0\},\quad r\ge1,
 \quad a_j\ge1.
\]
Let $A=\max_j a_j$, and let
\[
 d_\lambda
 =\#\{\alpha\in\Phi_G^+:\langle\lambda,\alpha^\vee\rangle>0\}.
\]
Then
\[
 r=\deg f+1\le d_\lambda+1\le|\Phi_G^+|+1\le\dim G+1,
\]
and P-F11 gives
\[
 p(f)<(d_\lambda+1)A+128(d_\lambda+1)^3A^3.
 \tag{P-F12a}
\]
In particular, for the plethysm and generalized plethysm problem of
GCT VI, Problem 1.1.2, Hypothesis 1.6.7 (PH3) implies Hypothesis 1.6.6
(PH2), with their stated parameter
$h=\dim G+\operatorname{ht}(\lambda)+\operatorname{ht}(\pi)$.
The same implication holds for connected subgroup restriction with
classical ambient $G$, using $\operatorname{rank}(x)$ as in GCT VI,
Section 3.4, and Hypotheses 3.3.3--3.3.4. This includes the usual
Kronecker restriction realization of Problem 1.1.1.

For $G=GL_m$, $m\ge1$, and a partition $\lambda$, the same positive form has a
binary description of length polynomial in
$m,A,\log(\lambda_1+1)$, with $\lambda_1=0$ for the empty partition.
Consequently PH3 also implies the existence of a polynomial-size shifted
polytope description through P-F2 in this expanded group-size parameter.
This last assertion is an existence and encoding bound, not a uniform
algorithm for constructing the positive form.
\end{proposition}

\begin{proof}
Complete reducibility gives
\[
 0\le f(n)\le\dim V_{n\lambda}(G).
\]
The Weyl dimension formula is
\[
 \dim V_{n\lambda}(G)
 =\prod_{\alpha\in\Phi_G^+}
   \frac{n\langle\lambda,\alpha^\vee\rangle
             +\langle\varrho_G,\alpha^\vee\rangle}
        {\langle\varrho_G,\alpha^\vee\rangle},
 \tag{P-F12b}
\]
where $\varrho_G$ is half the sum of positive roots. See
Borcherds--Haiman--Reshetikhin--Serganova--Johnson-Freyd,
\emph{Lectures in Lie Groups and Quantum Groups}, Part I (2016),
Theorem 6.1.1.30, printed pp.~102--103.
For reductive $G$, apply that formula to its semisimple derived Lie
algebra; the central character does not change the dimension.
Every nonconstant factor in (P-F12b) has positive slope, and their
number is $d_\lambda$. Hence $f(n)=O((n+1)^{d_\lambda})$.
The implied constant may depend on $\lambda$; its size is irrelevant to
this degree comparison.

We next bound the number of denominator factors without assuming that
the positive form came from a particular h.s.o.p. Choose $i_0$ with
$N_{i_0}>0$ and put $L=\operatorname{lcm}(a_1,\ldots,a_r)$.
For every nonnegative tuple $(y_1,\ldots,y_r)$ with $\sum_jy_j=k$,
the tuple $x_j=(L/a_j)y_j$ is integral and satisfies
$\sum_ja_jx_j=Lk$. Expanding the positive rational form therefore gives
\[
 f(Lk+i_0)\ge N_{i_0}\binom{k+r-1}{r-1}.
\]
Comparison with the Weyl bound forces $r-1\le d_\lambda$.
The possibly large integer $L$ is used only to prove this inequality;
it is not part of a required computational construction.

For completeness, expanding the same form shows
$f(n)=O((n+1)^{r-1})$: after fixing the numerator term and the first
$r-1$ nonnegative variables, the last variable is determined, and each
of the first $r-1$ variables is at most $n$ whenever the term contributes.
Together with the preceding lower bound on an infinite progression,
this proves $\deg f=r-1$. Finally
$|\Phi_G^+|=(\dim G-\operatorname{rank}G)/2\le\dim G$.
Global quasi-polynomiality gives
$\deg N<\sum_ja_j\le rA$ as in P-F9. Apply P-F11 and substitute
$r\le d_\lambda+1$ to obtain (P-F12a).

In Problem 1.1.2, $G=GL(V_\mu(H))$ and the inner weight $\mu$ is held
fixed while $\lambda$ and $\pi$ are stretched. Thus the preceding
multiplicity and the parameter $h$ are exactly the ones in
Hypotheses 1.6.6--1.6.7. If $A$ is bounded by a fixed polynomial in $h$,
then (P-F12a) bounds $p(f)$ by a fixed polynomial in $h$.
For a product of classical ambient groups, the numbers of positive
roots in types $A_\ell,B_\ell,C_\ell,D_\ell$ are respectively
$\ell(\ell+1)/2,\ell^2,\ell^2,\ell(\ell-1)$, and the sum for a product
is at most $(\operatorname{rank}G)^2$. The general linear central torus
does not add roots. Since the definition of $\operatorname{rank}(x)$
includes $\operatorname{rank}G$, one has
$r\le(\operatorname{rank}(x))^2+1$ in these cases.
The same argument proves the implication of Hypothesis 3.3.4 to 3.3.3.
For Kronecker coefficients use the restriction
$GL_a\times GL_b\to GL(\mathbb C^a\otimes\mathbb C^b)$, with the
third partition as the highest weight of the ambient general linear
group; this is the realization in Problem 1.1.1.
If all positive-degree multiplicities vanish, their constituent
polynomials are zero and their positivity index is zero. Such a ray
is handled separately: the isolated value $f(0)=1$ of the literal
branching formula is not part of a global quasi-polynomial and is
outside the hypothesis above.

It remains to prove the encoding assertion. Let $B=\deg N$ and write
$N(T)=\sum_{i=0}^B N_iT^i$. The reciprocal denominator has nonnegative
coefficients and constant coefficient one, so
\[
 N_i\le f(i)\le\dim V_{i\lambda}(GL_m)
 \le (i\lambda_1+m)^{m(m-1)/2}.
\]
The last inequality follows term by term from the $GL_m$ specialization
of (P-F12b): each factor is
$(i(\lambda_j-\lambda_k)+k-j)/(k-j)$ for $j<k$, and its denominator
is at least one. Since $B<rA$ and $r\le m(m-1)/2+1$, there are fewer
than $rA+1$ numerator entries, each requiring at most
\[
 O\!\left(1+m^2\log\bigl(2+rA\lambda_1+m\bigr)\right)
\]
bits. The denominator list requires $O(r\log(A+1))$ bits.
This is polynomial in $m,A,\log(\lambda_1+1)$.
Applying P-F2 to this explicit data gives the asserted polynomial-size
shifted description whenever such a form is supplied; existentially,
PH3 supplies a form with these bounds. The proof does not find that
form or certify its equality to the multiplicity series in polynomial
time.
\end{proof}

\paragraph{Scope and exact remaining gate.}
GCT VI itself remarks in Section 4.2.2 (printed p.~73) that PH2 may
follow from PH3. P-F11--P-F12 establish this implication for the stated
representation-theoretic families with an explicit bound. They do not
prove PH3. They also do not replace the geometric coefficient and basis
identifications needed to infer PH3 from a proposed purity construction.

The group-size parameter must be preserved. In generalized plethysm,
$\dim G=(\dim V_\mu(H))^2$ can be very large compared with the binary
encoding of $\mu$. GCT VI, Section 3.4, explicitly permits this compact
input; Hypothesis 1.6.4 requires its ambient polytope dimension to be
polynomial in the bitlength of $\mu$ and the heights of $\lambda,\pi$.
Thus the polynomial bounds just proved in $h$ do not automatically give
a uniform shifted model with that more compressed complexity bound.

No estimate $\deg N\le\deg f$ has been used. The numerator index in
the displays of GCT VI, Theorem 1.6.1(d) and equation (4.4), cannot
literally impose such a bound. Indeed the KM example has degree one,
so that bound would require a positive numerator $1+cT$ and two
denominator factors. Its value $\phi(1)=0$ forces $c=0$ and both
denominator weights to exceed one. Then $\phi(2)=\phi(3)=1$ forces
the weights to be two and three, but the resulting coefficient of
$T^4$ is one rather than $\phi(4)=2$, a contradiction. This identifies
an index clash in the printed display; the correct positive-form
consequence of Cohen--Macaulayness has an unrestricted numerator degree.
The required bound here is the proved properness inequality
$\deg N<\sum_j a_j$.
